\subsection{Indentación}

\subsubsection{Regla: cuatro espacios por nivel}

\textbf{Regla:} Se adopta la configuración predeterminada de formato de C\#: cada nivel de indentación usa cuatro espacios, el ancho de tabulación es cuatro y no se insertan caracteres de tabulación. La configuración equivalente es \texttt{indent\_style = space}, \texttt{indent\_size = 4} y \texttt{tab\_width = 4} (Microsoft, 2025e).

\textbf{Descripción:} Las convenciones de código de la documentación de .NET (Microsoft, 2025e), que la propia página ofrece para adoptarse tal cual o adaptarse a cada equipo, lo fijan así:
\begin{quote}
Use four spaces for indentation. Don't use tab characters. (sección Style guidelines, párr. 2)
\end{quote}
\noindent La misma página pide conservar la configuración predeterminada del editor, que indenta con cuatro caracteres y guarda los tabuladores como espacios (Microsoft, 2025e, sección Layout conventions, párr. 2).

\textbf{Código correcto:}
\begin{lstlisting}[style=csharpstyle]
public int RollDice()
{
    const int MinimumFaceValue = 1;
    const int MaximumFaceValue = 6;
    const int ExclusiveMaximumFaceValue = MaximumFaceValue + 1;
    int faceValue = Random.Shared.Next(MinimumFaceValue, ExclusiveMaximumFaceValue);

    if (faceValue == MaximumFaceValue)
    {
        return faceValue + BonusRoll();
    }

    return faceValue;
}
\end{lstlisting}

\textbf{Código incorrecto:}
\begin{lstlisting}[style=csharpstyle]
public int RollDice()
{
  const int MinimumFaceValue = 1;
  const int MaximumFaceValue = 6;
  const int ExclusiveMaximumFaceValue = MaximumFaceValue + 1;
  int faceValue = Random.Shared.Next(MinimumFaceValue, ExclusiveMaximumFaceValue);

  if (faceValue == MaximumFaceValue)
  {
      return faceValue + BonusRoll();
  }

  return faceValue;
}
\end{lstlisting}
